% GENERATED FILE. Edit canonical lessons, then run npm run generate:books.
% canonical-source-hash: fnv1a64:e1cd4b2c3253800d
% canonical-lessons: ES-C646-surtidor, ES-C646-repostar, ES-C646-deposito, ES-C646-combustible, ES-C646-kilometro

\chapter{Llenar el depósito}
\label{ch:es-llenar-el-deposito}
\hlchaptermodality{646}

\begin{chapteropening}
\textbf{By the end of this chapter:} \emph{I can fill up at a petrol pump and talk about the tank, fuel and kilometres.}

Complete the last lesson of chapter 646: I can fill up at a petrol pump and talk about the tank, fuel and kilometres.
\end{chapteropening}

\section[el surtidor]{el surtidor — the petrol pump; the jet (of a fountain)}

\label{lesson:ES-C646-surtidor}



\begin{quote}
Before the new one: say the Spanish for the tyre, then the Spanish for to get a puncture.
\end{quote}

\subsection*{You'll want to know: el surtidor}
\textbf{el surtidor} — \textquotedblleft{}the petrol pump; the jet (of a fountain)\textquotedblright{}.

\begin{quote}
Pare junto al surtidor número tres.
\end{quote}

\begin{quote}
\emph{Stop next to pump number three.}
\end{quote}

At a petrol station each pump is \emph{un surtidor}.

\begin{cousinweb}
\emph{surtidor} is built on \emph{surtir} \textquotedblleft{}to supply; to spout\textquotedblright{}. The origin of \emph{surtir} is not settled.
\end{cousinweb}

\begin{grammarlens}[title={What you've built}]
One more word for telling, explaining and describing at B1.
\end{grammarlens}

\subsection*{Your turn}
\begin{itemize}
\raggedright
  \item \emph{Say it:} \emph{el surtidor}
  \item \emph{Say it:} \emph{el surtidor}, once more
  \item \emph{Say it:} a sentence of your own with \emph{el surtidor}
  \item \emph{Recall:} say the Spanish for the tyre, then the Spanish for to get a puncture, then say \emph{el surtidor} again
\end{itemize}

\begin{tcolorbox}[breakable,colback=teal!4,colframe=teal!35!black,title={Before you move on}]
What does el surtidor mean? (\textquotedblleft{}The petrol pump; the jet (of a fountain)\textquotedblright{}.) Say it once more.
\end{tcolorbox}
\section[repostar]{repostar — to refuel, to fill up}

\label{lesson:ES-C646-repostar}



\begin{quote}
Before the new one: say the Spanish for to get a puncture, then the Spanish for the petrol pump.
\end{quote}

\subsection*{You'll want to know: repostar}
\textbf{repostar} — \textquotedblleft{}to refuel, to fill up\textquotedblright{}.

\begin{quote}
Tenemos que repostar antes de salir.
\end{quote}

\begin{quote}
\emph{We have to fill up before we leave.}
\end{quote}

A Spain word; in Latin America people say \emph{cargar gasolina} or \emph{echar gasolina}.

\begin{cousinweb}
\emph{repostar} is usually traced to \emph{reposte} \textquotedblleft{}store-room, larder\textquotedblright{}, from Latin \emph{repositum} \textquotedblleft{}something put away\textquotedblright{}, built on \emph{pōnere} \textquotedblleft{}to put\textquotedblright{}; it first meant taking on fresh supplies.
\end{cousinweb}

\begin{grammarlens}[title={What you've built}]
One more word for telling, explaining and describing at B1.
\end{grammarlens}

\subsection*{Your turn}
\begin{itemize}
\raggedright
  \item \emph{Say it:} \emph{repostar}
  \item \emph{Say it:} \emph{repostar}, once more
  \item \emph{Say it:} a sentence of your own with \emph{repostar}
  \item \emph{Recall:} say the Spanish for to get a puncture, then the Spanish for the petrol pump, then say \emph{repostar} again
\end{itemize}

\begin{tcolorbox}[breakable,colback=teal!4,colframe=teal!35!black,title={Before you move on}]
What does repostar mean? (\textquotedblleft{}To refuel, to fill up\textquotedblright{}.) Say it once more.
\end{tcolorbox}
\section[el depósito]{el depósito — the tank; the deposit}

\label{lesson:ES-C646-deposito}



\begin{quote}
Before the new one: say the Spanish for the petrol pump, then the Spanish for to refuel.
\end{quote}

\subsection*{You'll want to know: el depósito}
\textbf{el depósito} — \textquotedblleft{}the tank; the deposit\textquotedblright{}.

\begin{quote}
Llene el depósito, por favor.
\end{quote}

\begin{quote}
\emph{Fill the tank, please.}
\end{quote}

At a bank \emph{el depósito} is a deposit.

\begin{cousinweb}
\emph{depósito} is a learned borrowing from Latin \emph{dēpositum} \textquotedblleft{}something laid down, entrusted\textquotedblright{}, from \emph{dēpōnere} \textquotedblleft{}to put down\textquotedblright{} (Lewis \& Short), built on \emph{pōnere} \textquotedblleft{}to put\textquotedblright{}.
\end{cousinweb}

\begin{grammarlens}[title={What you've built}]
One more word for telling, explaining and describing at B1.
\end{grammarlens}

\subsection*{Your turn}
\begin{itemize}
\raggedright
  \item \emph{Say it:} \emph{el depósito}
  \item \emph{Say it:} \emph{el depósito}, once more
  \item \emph{Say it:} a sentence of your own with \emph{el depósito}
  \item \emph{Recall:} say the Spanish for the petrol pump, then the Spanish for to refuel, then say \emph{el depósito} again
\end{itemize}

\begin{tcolorbox}[breakable,colback=teal!4,colframe=teal!35!black,title={Before you move on}]
What does el depósito mean? (\textquotedblleft{}The tank; the deposit\textquotedblright{}.) Say it once more.
\end{tcolorbox}
\section[el combustible]{el combustible — the fuel}

\label{lesson:ES-C646-combustible}



\begin{quote}
Before the new one: say the Spanish for to refuel, then the Spanish for the tank.
\end{quote}

\subsection*{You'll want to know: el combustible}
\textbf{el combustible} — \textquotedblleft{}the fuel\textquotedblright{}.

\begin{quote}
El combustible está muy caro.
\end{quote}

\begin{quote}
\emph{Fuel is very expensive.}
\end{quote}

As an adjective it means something that can burn.

\begin{cousinweb}
\emph{combustible} is a learned borrowing built on Latin \emph{combūrere} \textquotedblleft{}to burn up, to consume\textquotedblright{} (Lewis \& Short), through its past form \emph{combustus}.
\end{cousinweb}

\begin{grammarlens}[title={What you've built}]
One more word for telling, explaining and describing at B1.
\end{grammarlens}

\subsection*{Your turn}
\begin{itemize}
\raggedright
  \item \emph{Say it:} \emph{el combustible}
  \item \emph{Say it:} \emph{el combustible}, once more
  \item \emph{Say it:} a sentence of your own with \emph{el combustible}
  \item \emph{Recall:} say the Spanish for to refuel, then the Spanish for the tank, then say \emph{el combustible} again
\end{itemize}

\begin{tcolorbox}[breakable,colback=teal!4,colframe=teal!35!black,title={Before you move on}]
What does el combustible mean? (\textquotedblleft{}The fuel\textquotedblright{}.) Say it once more.
\end{tcolorbox}
\section[el kilómetro]{el kilómetro — the kilometre}

\label{lesson:ES-C646-kilometro}



\begin{quote}
Before the new one: say the Spanish for the tank, then the Spanish for the fuel.
\end{quote}

\subsection*{You'll want to know: el kilómetro}
\textbf{el kilómetro} — \textquotedblleft{}the kilometre\textquotedblright{}.

\begin{quote}
Faltan diez kilómetros.
\end{quote}

\begin{quote}
\emph{There are ten kilometres to go.}
\end{quote}

Often written \emph{km}. \emph{El kilómetro cero} is the point distances are measured from.

\begin{cousinweb}
\emph{kilómetro} was borrowed from French \emph{kilomètre}, a word made from Greek \emph{khílioi} \textquotedblleft{}a thousand\textquotedblright{} and \emph{métron} \textquotedblleft{}measure\textquotedblright{}.
\end{cousinweb}

\begin{grammarlens}[title={What you've built}]
One more word for telling, explaining and describing at B1.
\end{grammarlens}

\subsection*{Your turn}
\begin{itemize}
\raggedright
  \item \emph{Say it:} \emph{el kilómetro}
  \item \emph{Say it:} \emph{el kilómetro}, once more
  \item \emph{Say it:} a sentence of your own with \emph{el kilómetro}
  \item \emph{Recall:} say the Spanish for the tank, then the Spanish for the fuel, then say \emph{el kilómetro} again
\end{itemize}

\begin{tcolorbox}[breakable,colback=teal!4,colframe=teal!35!black,title={Before you move on}]
What does el kilómetro mean? (\textquotedblleft{}The kilometre\textquotedblright{}.) Say it once more.
\end{tcolorbox}
